% Generated by `taskgen course examples` from dlp.rho@composite (seed 20).
% Do not edit: change examples.yaml and regenerate.
\begin{example}[$\rho$-алгоритм, $\gcd(\beta - \beta', n) > 1$]\label{ex:rho-composite}
Решить задачу дискретного логарифмирования $5^x \equiv 29\ (\textrm{mod}\ 37)$ $\rho$-алгоритмом Полларда с поиском цикла по Флойду; $\ord(5) = 36 = 2^{2} \cdot 3^{2}$. Последовательность $x_{i+1} = f(x_i)$, $x_i = 5^{\alpha_i} \cdot 29^{\beta_i} \bmod 37$ (показатели по модулю $36$), где\[f(x) = \begin{cases} x^2, & x \equiv 0,\\ 29\, x, & x \equiv 1,\\ 5\, x, & x \equiv 2 \end{cases} \pmod 3, \qquad (\alpha, \beta) \mapsto \begin{cases} (2\alpha, 2\beta),\\ (\alpha, \beta + 1),\\ (\alpha + 1, \beta), \end{cases}\]начиная с $x_0 = 5^{20} \cdot 29^{26} \equiv 28$ ($\alpha_0 = 20$, $\beta_0 = 26$). Выписать таблицу $(x_i, \alpha_i, \beta_i)$, $(x_{2i}, \alpha_{2i}, \beta_{2i})$ до совпадения $x_i = x_{2i}$ и найти $x$.

\par\noindent\textit{Решение.} Вычисляем «черепаху» $x_i = f(x_{i-1})$ и «зайца» $x_{2i} = f(f(x_{2i-2}))$:\[\begin{array}{r|rrr|rrr}i & x_i & \alpha_i & \beta_i & x_{2i} & \alpha_{2i} & \beta_{2i} \\ \hline 0 & 28 & 20 & 26 & 28 & 20 & 26 \\ 1 & 35 & 20 & 27 & 27 & 21 & 27 \\ 2 & 27 & 21 & 27 & 19 & 7 & 18 \\ 3 & 26 & 6 & 18 & 16 & 14 & 2 \\ 4 & 19 & 7 & 18 & 26 & 15 & 3 \\ 5 & 33 & 7 & 19 & 33 & 16 & 4 \\ \end{array}\]Совпадение при $i = 5$: $x_{5} = x_{10} = 33$ (у орбиты $x_0$ предпериод $T = 3$ и период $M = 5$; $i = 5$ --- наименьшее кратное $M$, не меньшее $T$).\par Из $g^{7} h^{19} = g^{16} h^{4}$ и $h = g^x$ получаем $(19 - 4)\, x \equiv 16 - 7\ (\textrm{mod}\ 36)$, т.\,е. $15\, x \equiv 9\ (\textrm{mod}\ 36)$. $d = \gcd(15, 36) = 3 > 1$ и $d \mid 9$: сокращаем на $d$: $5\, x \equiv 3\ (\textrm{mod}\ 12)$; по модулю $36$ получится $d = 3$ кандидатов.  Обратный к $5$ по модулю $12$ (расширенный алгоритм Евклида для $(12, 5)$):
\[\begin{array}{r|r|r|r}
r_i & q_i & u_i & v_i \\ \hline
12 &  & 1 & 0 \\
5 &  & 0 & 1 \\
2 & 2 & 1 & -2 \\
1 & 2 & -2 & 5 \\
0 & 2 & 5 & -12 \\
\end{array}\qquad r_i = u_i \cdot 12 + v_i \cdot 5\]

$5^{-1} \equiv 5\ (\textrm{mod}\ 12)$.
\par $\tilde x \equiv 3 \cdot 5 \equiv 3\ (\textrm{mod}\ 12)$. Кандидаты $x = \tilde x + k \cdot 12$, $k = 0, \dots, 2$; проверяем $g^x \equiv h$: $5^{3} \equiv 14$, $5^{15} \equiv 29$\ (\checkmark), $5^{27} \equiv 31$.\[x = 15.\]
\end{example}
